\begin{derivation}[从\eqnrefs{eq:eeff-chiral-coeff-dp11}到\eqnrefs{eq:eeff-spin-summed-dp11}]
考虑无质量过程 $e^-(p_1)e^+(p_2)\to f(k_1)\bar f(k_2)$，固定电子流手征 $i=L/R$ 与末态流手征 $j=L/R$。相应的流乘积为

\begin{equation*}
\mathcal J_{ij}
=[\bar v(p_2)\gamma^\mu P_i u(p_1)]
[\bar u(k_1)\gamma_\mu P_j v(k_2)].
\end{equation*}
对末态自旋求和并对初态自旋平均后，基本张量为
\begin{equation*}
\frac14\Sigma_{ij}
=\frac14\Tr(\slashed p_2\gamma^\mu P_i\slashed p_1\gamma^\nu P_i)
\Tr(\slashed k_1\gamma_\mu P_j\slashed k_2\gamma_\nu P_j).
\end{equation*}
利用
$P_L\slashed p=\slashed pP_R$、$P_R\slashed p=\slashed pP_L$ 以及
$P_R\gamma^\nu=\gamma^\nu P_L$，左手迹中的两个投影满足
\begin{align*}
 \Tr(\slashed a\gamma^\mu P_L\slashed b\gamma^\nu P_L)
 &=\Tr(\slashed a\gamma^\mu\slashed bP_R\gamma^\nu P_L),\\
 &=\Tr(\slashed a\gamma^\mu\slashed b\gamma^\nu P_L^2),\\
 &=\Tr(\slashed a\gamma^\mu\slashed b\gamma^\nu P_L).
\end{align*}
右手迹以 $L\leftrightarrow R$ 进行同样的投影代数。以左手投影为例，
\begin{align*}
\Tr(\slashed a\gamma^\mu\slashed b\gamma^\nu P_L)
&=\frac12\Tr(\slashed a\gamma^\mu\slashed b\gamma^\nu)
-\frac12\Tr(\slashed a\gamma^\mu\slashed b\gamma^\nu\gamma^5)\\
&=2\left[S^{\mu\nu}(a,b)+\ii A^{\mu\nu}(a,b)\right],
\end{align*}
其中
\begin{align*}
S^{\mu\nu}(a,b)
&\equiv a^\mu b^\nu+a^\nu b^\mu-\eta^{\mu\nu}a\cdot b,\\
A^{\mu\nu}(a,b)
&\equiv\epsilon^{\alpha\mu\beta\nu}a_\alpha b_\beta.
\end{align*}
右手投影只把 $A^{\mu\nu}$ 的符号反转。现在把“手征选择”真正做成张量收缩。由于一个张量对称、另一个反对称，
\begin{equation*}
S^{\mu\nu}(a,b)A_{\mu\nu}(c,d)=0.
\end{equation*}
对称部分直接展开为
\begin{equation*}
S^{\mu\nu}(a,b)S_{\mu\nu}(c,d)
=2\left[(a\cdot c)(b\cdot d)+(a\cdot d)(b\cdot c)\right].
\end{equation*}
反对称部分使用
\begin{equation*}
\epsilon^{\alpha\mu\beta\nu}\epsilon_{\gamma\mu\delta\nu}
=-2\left(\delta^\alpha_{\gamma}\delta^\beta_{\delta}
-\delta^\alpha_{\delta}\delta^\beta_{\gamma}\right)
\end{equation*}
得到
\begin{equation*}
(\ii A^{\mu\nu})(\ii A_{\mu\nu})
=2\left[(a\cdot d)(b\cdot c)-(a\cdot c)(b\cdot d)\right].
\end{equation*}
因此同手征时对称部分与反对称部分相加，只保留第二种配对；异手征时其中一个 $A$ 反号，只保留第一种配对。取
$(a,b,c,d)=(p_2,p_1,k_1,k_2)$，并记住每个手征迹前的因子 $2$ 以及初态平均 $1/4$，便有
\begin{align*}
\frac14\Sigma_{LL}
=\frac14\Sigma_{RR}
&=4(p_1\cdot k_1)(p_2\cdot k_2),\\
\frac14\Sigma_{LR}
=\frac14\Sigma_{RL}
&=4(p_1\cdot k_2)(p_2\cdot k_1).
\end{align*}

在质心系令 $x=\cos\theta$，无质量运动学给出
\begin{align*}
2p_1\cdot k_1&=-t=\frac{s}{2}(1-x),\\
2p_1\cdot k_2&=-u=\frac{s}{2}(1+x),
\end{align*}
并且由质心系的交换对称性
\begin{equation*}
p_2\cdot k_2=p_1\cdot k_1,
\qquad
p_2\cdot k_1=p_1\cdot k_2.
\end{equation*}
于是两类手征块分别产生 $(1-x)^2$ 与 $(1+x)^2$。按照 $C_{LL},C_{LR},C_{RL},C_{RR}$ 的外流手征约定重新排列后，即得到\eqnrefs{eq:eeff-spin-summed-dp11}。因此二次角权重由手征投影下的对称四伽马迹与反对称 Levi--Civita 迹共同确定。
\end{derivation}

% $A_{\rm FB}$ 的 $3/4$ 因子只需两个二次多项式的半区间积分，正文保留结论，不另设推导框。

